\subsection{A horizontal filling that removes the resolved-tail period}

\begin{proposition}[AT19: Proved; an actual zero-period lift for the quadratic sphere disk]
\label{AT19}
Let $R=\mathbb Z_{(2)}$, and let
\[
\begin{gathered}
 X=G_R\otimes_R\operatorname{FreePD}_R
          \bigl(x_0K(2)\xrightarrow{\,d\,}y_0K(2)\bigr),\\
 |x_0|=1,\qquad |y_0|=2,\qquad
 \operatorname{wt}(x_0)=\operatorname{wt}(y_0)=b\geq3.
\end{gathered}
 \tag{AT.105}
\]
Here the coefficient action is that of $G_R$, and the tensor
product denotes its free relative extension.
Suppose actual strict maps $f_A:X\to A$, $f_B:X\to B$ are
given, with torsionfree admissible targets. Put $s=2b$ and
assume
\[
 b\geq\max\{\mu_A(1),\mu_B(1)\},\qquad
 s+1\geq\max\{\lambda_A(4),\lambda_B(4)\}.
 \tag{AT.106}
\]
Use the coefficient resolution and the resolved-tail algebra
of A-S09-22, with $L=I_1\langle s\rangle$ and twist
$z_0=\gamma_2(y_0)$, $\eta_0=x_0y_0$.

The maps $f_A,f_B$ have actual extensions to that resolved-tail
algebra for which the corrected pair
$(\zeta_{\mathrm{tw}},B)$ of (AT.90) admits simultaneous
relative two-cell target lifts:
\[
 \begin{aligned}
 dU_A&=f_A(\zeta_{\mathrm{tw}}),&
 dV_A&=2U_A-f_A(B),\\
 U_A&\in\operatorname{Hom}(K(2),A^2),&
 V_A&\in\operatorname{Hom}(K(2),A^1),
 \end{aligned}
 \tag{AT.107}
\]
and the analogous equations in $B$.
The elements have weight $2s+5$.
Thus the actual relative repair (AT.92) extends both maps.
When the AT18 period comparison is used, its residue is zero.
No vertical comparison at the original weight $s$ is needed.

A-SAT01 at cutoff $s+1$ constructs an actual iterated
saturation of this resolved algebra and extending maps
to both targets. The same relative repair extends those
maps as well: the elements in (AT.107) continue to satisfy
the very same generator equations.
This assertion does not supply eventual admissibility of
the repaired source.
\end{proposition}

\begin{proof}
We construct the values for $A$; the construction for $B$
is independent. Let
\[
                  \xi=f_A(x_0)\in N_2(A^{b,1}).
\]
All its faces are zero, since $x_0$ is the sphere generator.
The row-one bound in (AT.106) therefore supplies
\[
             \alpha\in N_3(A^{b,1}),\qquad d_0\alpha=\xi.
 \tag{AT.108}
\]
Write $D(3)$ for the Dold--Kan image of the normalized
disk $R[3]\xrightarrow{1}R[2]$; its bottom is $K(2)$.
Equation (AT.108) gives an actual linear simplicial map
$D(3)\to A^1$ extending the old map on that bottom.

We specify the nonlinear expressions by normalized chain
maps, not by declaring a quadratic expression to be linear
on arbitrary sums of parameters. In normalized degree three
take
\[
              \widetilde z_{\rm top}=\gamma_2(d\alpha),
                 \qquad
              \widetilde\eta_{\rm top}=\alpha\,d\alpha.
 \tag{AT.109}
\]
Every face $d_i$, $i>0$, is zero. Their remaining faces are
$\gamma_2(d\xi)$ and $\xi\,d\xi$, respectively.
Those are normalized degree-two cycles, since every face
of $\xi$ is zero. Hence the assignments
\[
 \begin{array}{c|cc}
 &R[3]&R[2]\\\hline
 \widetilde z&\gamma_2(d\alpha)&\gamma_2(d\xi)\\
 \widetilde\eta&\alpha\,d\alpha&\xi\,d\xi
 \end{array}
 \tag{AT.110}
\]
are actual maps of normalized chain complexes from the disk.
Dold--Kan gives actual linear simplicial maps on $D(3)$.
Their vertical identity $d\widetilde\eta=2\widetilde z$
holds on both normalized generators by the strict PD rules;
it therefore holds everywhere on these linear maps.
Coefficient multiplication extends them to the required
$I_1\langle s\rangle$ parameter. This proves all parameter
compatibilities without an additivity assumption on $\gamma_2$.

Apply the full free coefficient complex $P$ from AT16,
with this disk-parameter twist. Put
\[
          \mathcal D_A=\operatorname{Hom}_{sR\text{-}\mathrm{Mod}}
                                      (D(3),A)=N_3(A).
 \tag{AT.111}
\]
The equality is an actual cochain-module identification:
a normalized chain map out of the disk is determined by
the image of its degree-three generator.
The disk is projective in simplicial $R$-modules, so this
Hom functor is exact and commutes with vertical cohomology.
In any weight at least $\lambda_A(4)$, its cohomology
through form degree four is zero. Positive degrees vanish
by admissibility; in degree zero the coefficient comparison
is the constant simplicial module $S_R^w$, whose third
normalization is zero.

Every free coefficient generator of $P$ has weight at
least $s+1$. The second bound in (AT.106) therefore applies
to every group on the finite total-degree-four diagonal
used in AT16. With the disk target (AT.111), all entries
on that diagonal vanish. Consequently the twisting
degree-one Hom cocycle has a degree-zero primitive.
Equivalently, applying the free resolution
$F_r=P^{3-r}$ of $L/2L$, the first-quadrant Hom comparison
in total degree four has zero target cohomology in every
degree $0,\ldots,4$. This proves the required vanishing
without a cohomology assumption at weight $s$.

The terminal equation of this primitive is
$h_0d_P^{-1}=0$. It makes $h_0$ descend to the actual
bottom quotient $W^0$, exactly as in AT16.
We obtain actual $S_R$-linear generator values
\[
                    h:W\otimes D(3)\longrightarrow A
 \tag{AT.112}
\]
with their prescribed affine differential.
Restriction to the bottom $K(2)$ gives the required
strict map from the resolved-tail algebra to $A$.
Torsionfreeness of $A$ makes it factor through the
specified torsionfree quotient.

Here is also a direct algebraic description of the
one-simplex-higher values, useful for checking all
polynomial expressions. In simplex two the source
algebra is freely specified over $(G_R)_2$ by the
old disk generator and the resolved $W$ generators,
with the displayed relations. Send $(G_R)_2$ by its
actual degeneracy $s_0$ followed by $f_A$ into $A_3$;
send $x_0$ to $\alpha$, $y_0$ to $d\alpha$; and send
each $W$ generator to the normalized top value in
(AT.112). The checked differential and coefficient
relations give an actual strict PD cochain map
\[
                         \widetilde f:Y_2^{\rm tf}\longrightarrow A_3.
 \tag{AT.113}
\]
Its $d_0$ face is the chosen target map on $Y_2^{\rm tf}$.
For every $i>0$, its face kills all old disk and resolved
generators. In particular every polynomial expression
in the ideal generated by those variables has zero
$i$th face after (AT.113). This statement uses actual
generator maps; it does not assert a linearization of a
polynomial function on an arbitrary parameter.

Apply (AT.113) to the two expressions $B$ and
$\zeta_{\rm tw}$. Both are in that ideal.
Their top values are normalized three-simplices,
their $d_0$ values are the corresponding prescribed
normalized two-simplices, and their other faces vanish.
The source identity $dB=2\zeta_{\rm tw}$ is preserved.
They therefore give actual linear disk maps
\[
 \widetilde B:D(3)\to A^{2s+5,2},\qquad
 \widetilde\zeta:D(3)\to A^{2s+5,3},
 \qquad d\widetilde B=2\widetilde\zeta,
 \tag{AT.114}
\]
again by the two normalized generators, with
$d\widetilde\zeta=0$.

This weight is at least $s+1$, so (AT.111) has zero
cohomology in degrees three and two there.
Choose actual elements in that Hom complex satisfying
\[
 d\widetilde U=\widetilde\zeta,\qquad
 d\widetilde V=2\widetilde U-\widetilde B.
 \tag{AT.115}
\]
The second right-hand side is closed by (AT.114).
Restrict (AT.115) to the bottom $K(2)$.
This gives precisely $U_A,V_A$ in (AT.107).
Thus the free relative PD universal property supplies
the actual extended target map after the repair.
This construction proves zero residue directly by
integral primitives; alternatively reduction of
(AT.115) makes $f_A(B)$ a boundary modulo two in
the relative complex.

Repeat all steps for $B$ using its own row filling
and its own Hom primitives. The source formulas and
the repaired source are the same for both maps.
Finally apply A-SAT01 to the actual torsionfree resolved
algebra, with these chosen maps and cutoff $s+1$.
Its simplex-zero algebra is $S_R$, and its coefficient
$H^0$ and rational coefficient comparison were proved
by the explicit AT16 contraction and target-integrality
argument. Those arguments apply to the actual maps
just constructed. The target degree-zero and degree-one
comparisons above $s+1$ follow from (AT.106).
Thus every hypothesis of A-SAT01 holds and it supplies
the actual saturation and extending target maps.
They retain the same values on $B$ and $\zeta_{\rm tw}$.
The two new generator assignments (AT.107) therefore
extend those maps as well. This proves the asserted
scope of the construction.
\end{proof}
